% ================================================================
% PRESENTATION BEAMER — Battle Météorage 2026
% GBS Survival Analysis — basée sur le code réel
% ================================================================
\documentclass[aspectratio=169,t,11pt]{beamer}

\usepackage[T1]{fontenc}
\usepackage[utf8]{inputenc}
\usepackage[french]{babel}
\usepackage{lmodern}
\usepackage{amsmath,amssymb}
\usepackage{booktabs}
\usepackage{array}
\usepackage{tikz}
\usetikzlibrary{arrows.meta,decorations.pathreplacing,calc}
\usepackage[most]{tcolorbox}
\usepackage{listings}
\tcbuselibrary{listings}
\usepackage{xcolor}
\usepackage{graphicx}
\graphicspath{{notebooks/}{results/}{_archive/}}

% ── Couleurs ────────────────────────────────────────────────────
\definecolor{metBlue}{RGB}{14,52,100}
\definecolor{metMid}{RGB}{46,117,182}
\definecolor{metLight}{RGB}{191,215,237}
\definecolor{metGreen}{RGB}{10,150,100}
\definecolor{metOrange}{RGB}{210,110,0}
\definecolor{metRed}{RGB}{190,45,45}
\definecolor{codebg}{RGB}{245,247,250}
\definecolor{codeframe}{RGB}{170,185,210}

% ── Thème Beamer ────────────────────────────────────────────────
\usetheme{Boadilla}
\setbeamertemplate{navigation symbols}{}
\setbeamercolor{structure}{fg=metBlue}
\setbeamercolor{frametitle}{bg=metBlue,fg=white}
\setbeamercolor{block title}{bg=metMid,fg=white}
\setbeamercolor{block body}{bg=metLight!25}
\setbeamercolor{block title alerted}{bg=metRed,fg=white}
\setbeamercolor{block body alerted}{bg=metRed!8}
\setbeamercolor{block title example}{bg=metGreen,fg=white}
\setbeamercolor{block body example}{bg=metGreen!10}
\setbeamerfont{frametitle}{size=\large,series=\bfseries}

\setbeamertemplate{headline}{%
  \begin{beamercolorbox}[wd=\paperwidth,ht=2pt]{palette primary}\end{beamercolorbox}}
\setbeamertemplate{footline}{%
  \begin{beamercolorbox}[wd=\paperwidth,ht=15pt,dp=4pt,
    leftskip=8pt,rightskip=8pt]{palette secondary}%
    \small\color{white}Battle Météorage 2026 · GBS Survival Analysis%
    \hfill\insertframenumber/\inserttotalframenumber\hspace{4pt}%
  \end{beamercolorbox}}

\setbeamersize{text margin left=7pt,text margin right=7pt}
\addtobeamertemplate{block begin}{}{\vspace{-2pt}}
\addtobeamertemplate{block end}{\vspace{-2pt}}{}

% ── tcolorbox ───────────────────────────────────────────────────
\tcbset{mybox/.style={
  arc=3pt,boxrule=1pt,
  top=3pt,bottom=3pt,left=5pt,right=5pt,
  fonttitle=\small\bfseries}}

% ── Style Python pour listings ──────────────────────────────────
\lstdefinestyle{pycode}{
  language=Python,
  basicstyle=\ttfamily\scriptsize,
  keywordstyle=\bfseries\color{metBlue},
  stringstyle=\color{metGreen!70!black},
  commentstyle=\itshape\color{gray!80},
  showstringspaces=false,
  breaklines=true,
  tabsize=2,
  upquote=true,
}

% tcblisting pour les blocs de code dans les slides
\newtcblisting{codeslide}[1]{
  arc=3pt, boxrule=0.8pt,
  top=3pt, bottom=3pt, left=5pt, right=5pt,
  colback=codebg, colframe=codeframe,
  title={#1},
  fonttitle=\scriptsize\bfseries\sffamily\color{metBlue},
  listing only,
  listing options={style=pycode},
}

% ================================================================
\begin{document}

% ── TITRE ───────────────────────────────────────────────────────
{
\setbeamertemplate{footline}{}
\setbeamertemplate{headline}{}
\begin{frame}[plain]
\begin{tikzpicture}[remember picture,overlay]
  \fill[metBlue] (current page.south west) rectangle (current page.north east);
  \fill[metGreen] ([yshift=-1pt]current page.north west)
    rectangle ([yshift=5pt]current page.north east);
  \fill[metMid!50] (current page.south west)
    rectangle ([yshift=2.6cm]current page.south east);
  \node[font=\LARGE\bfseries,text=white,anchor=center]
    at ([yshift=2.1cm]current page.center)
    {Gradient Boosting Survival Analysis};
  \node[font=\large,text=white,anchor=center]
    at ([yshift=1.2cm]current page.center)
    {Prédiction de fin d'alertes foudre -- Battle Météorage 2026};
  \draw[metGreen,line width=2pt]
    ([xshift=3cm]current page.west|-current page.center)
    -- ([xshift=-3cm]current page.east|-current page.center);
  \node[font=\normalsize,text=metLight,anchor=center]
    at ([yshift=-0.2cm]current page.center)
    {Modèle · Features · Code réel · Résultats};
\end{tikzpicture}
\end{frame}
}

% ================================================================
% 1. VARIABLE CIBLE
% ================================================================
\begin{frame}[fragile,shrink=20]{Variable cible : le silence inter-éclairs}
\vspace{-2pt}
\begin{columns}[T]

\column{0.52\textwidth}
\begin{codeslide}{build\_survival() -- modelisation\_v7.ipynb}
def build_survival(df, max_gap=30):
    g = df.groupby(['airport','airport_alert_id'])
    df['next_date'] = g['date'].shift(-1)
    df['duree'] = ((df['next_date'] - df['date'])
                   .dt.total_seconds() / 60)
    df['event'] = df['duree'].notna() & (df['duree']<=max_gap)
    df.loc[~df['event'], 'duree'] = max_gap  # censure a 30
    return df
# --> Events : 53 972 / 56 599  (95.4%)
# --> Censure :  2 627 derniers eclairs (4.6%)
\end{codeslide}
\vspace{3pt}
\begin{tcolorbox}[mybox,colback=metGreen!10,colframe=metGreen,
  title={Split temporel}]
\scriptsize
\textbf{Train 2016-2020 :} 41\,509 éclairs\\
\textbf{Test  2021-2022 :} 15\,090 éclairs\\
\textbf{Jury  2023-2025 :} 80\,186 éclairs (jamais vus)
\end{tcolorbox}

\column{0.46\textwidth}
\begin{tcolorbox}[mybox,colback=metMid!10,colframe=metMid,
  title={Variable $T_i$ : silence inter-éclairs}]
\small
$T_i$ = délai (min) jusqu'au prochain CG dans la même alerte,
plafonné à 30 min.\\[4pt]
\textbf{Event} $\delta_i=1$ : un éclair suit $\le$ 30 min.\\
\textbf{Censure} $\delta_i=0$ : alerte fermée, $T_i=30$.
\end{tcolorbox}
\vspace{4pt}
\begin{exampleblock}{Pourquoi la censure change tout}
\small
Ignorer les censurés biaise $\hat{T}$ vers le bas.\\
L'analyse de survie (Cox) les inclut correctement
via l'ensemble à risque $R(T_i)$.
\end{exampleblock}
\end{columns}
\end{frame}

% ================================================================
% 2. FEATURES
% ================================================================
\begin{frame}[fragile,shrink=25]{Feature engineering -- 13 variables construites}
\vspace{-2pt}
\begin{columns}[T]

\column{0.50\textwidth}
\begin{codeslide}{build\_features\_v7() -- extrait modelisation\_v7.ipynb}
# Cycliques : arbres ne voient pas que 23h ~ 1h
h = df['date'].dt.hour + df['date'].dt.minute/60
df['h_cos'] = np.cos(2*np.pi*h/24)
df['h_sin'] = np.sin(2*np.pi*h/24)
df['doy_cos'] = np.cos(2*np.pi*doy/365)
df['doy_sin'] = np.sin(2*np.pi*doy/365)
df['saison'] = ((df['date'].dt.month%12)//3)+1
# Cadence
df['silence_min'] = ((df['date']-prev)
    .dt.total_seconds()/60).fillna(30).clip(0,60)
df['freq_5min'] = (1/df['silence_min']
    .clip(lower=0.5)).clip(upper=10)
# Distance
df['dist_avg_5'] = g['dist'].transform(
    lambda x: x.rolling(5,min_periods=1).mean())
df['dist_min_so_far'] = g['dist'].cummin()
# v7 : LabelEncoder fitte sur TRAIN uniquement
df['airport_enc'] = le.transform(df['airport'])
\end{codeslide}

\column{0.48\textwidth}
\begin{tcolorbox}[mybox,colback=metBlue!5,colframe=metBlue,
  title={13 features -- 4 familles}]
\scriptsize
\textbf{Temporelles cycliques (5) :}\\
\texttt{h\_cos, h\_sin, doy\_cos, doy\_sin, saison}\\[3pt]
\textbf{Cadence (2) :}\\
\texttt{silence\_min, freq\_5min}\\[3pt]
\textbf{Distance (3) :}\\
\texttt{dist\_centre, dist\_avg\_5, dist\_min\_so\_far}\\[3pt]
\textbf{Alerte / aéroport (3) :}\\
\texttt{rang, rang\_norm, airport\_enc}
\end{tcolorbox}
\vspace{3pt}
\begin{tcolorbox}[mybox,colback=metGreen!10,colframe=metGreen,
  title={Sortie notebook}]
\scriptsize
\texttt{X\_train : (41509, 13)} \quad \texttt{X\_test : (15090, 13)}\\
\texttt{StandardScaler fitté sur train uniquement.}\\
\texttt{LabelEncoder : ['Ajaccio',..'Pise'] -> [0..4]}
\end{tcolorbox}
\end{columns}
\end{frame}

% ================================================================
% 3. MODELE ET ALGO
% ================================================================
\begin{frame}[shrink=10]{Le modèle : Cox à risques proportionnels + Boosting}
\vspace{-2pt}
\begin{columns}[T]

\column{0.50\textwidth}
\begin{block}{Hypothèse Cox (PH)}
\footnotesize
\[ \lambda(t \mid x_i) = \lambda_0(t) \cdot e^{f(x_i)} \]
$\lambda_0(t)$ : risque de base non spécifié.
$f(x_i)$ : score individuel appris.\\[3pt]
Survie individuelle via estimateur de Breslow :
\[ \hat{S}(t \mid x_i) = \exp\!\bigl(-\hat{\Lambda}_0(t)\,e^{\hat{f}(x_i)}\bigr) \]
\end{block}
\vspace{3pt}
\begin{block}{Vraisemblance partielle (perte du boosting)}
\footnotesize
\[
\mathcal{L}(f) = \sum_{i:\delta_i=1}
\!\Bigl[f(x_i) - \log\!\!\sum_{j \in R(T_i)}\!\!e^{f(x_j)}\Bigr]
\]
Censure incluse via l'ensemble à risque $R(T_i)$.
\end{block}

\column{0.48\textwidth}
\begin{tcolorbox}[mybox,colback=metBlue!5,colframe=metBlue,
  title={Boucle de boosting}]
\footnotesize
\textbf{Init} : $f_0 = 0$.\\[2pt]
\textbf{Iter $m$} : calculer le pseudo-résidu
\[r_i^{(m)} = -\frac{\partial\mathcal{L}}{\partial f(x_i)}\Big|_{f_{m-1}}\]
Ajuster un arbre $h_m$ sur $(x_i, r_i^{(m)})$.
Mettre à jour $f_m = f_{m-1} + \nu \cdot h_m$.\\[2pt]
\textbf{Sortie} : $\hat{f} = f_M$, puis Breslow pour $\hat{\Lambda}_0$.
\end{tcolorbox}
\vspace{3pt}
\begin{exampleblock}{Pourquoi boosting sur Cox ?}
\footnotesize
Cox linéaire ignore les interactions (saison $\times$ aéroport,
silence $\times$ distance). Le boosting apprend $f(x_i)$
non-linéaire sans imposer de forme sur $\lambda_0(t)$.
\end{exampleblock}
\end{columns}
\end{frame}

% ================================================================
% 4. ENTRAINEMENT
% ================================================================
\begin{frame}[fragile,shrink=20]{Entraînement du modèle}
\vspace{-2pt}
\begin{columns}[T]

\column{0.50\textwidth}
\begin{codeslide}{modelisation\_v7.ipynb}
from sksurv.ensemble import (
    GradientBoostingSurvivalAnalysis)
from sksurv.util import Surv

y_train = Surv.from_arrays(
    df_train['event'], df_train['duree'])

gbs_v7 = GradientBoostingSurvivalAnalysis(
    n_estimators=50, learning_rate=0.15,
    max_depth=3, subsample=0.9,
    min_samples_leaf=10,
    loss='coxph', random_state=42,
)
gbs_v7.fit(X_train_s, y_train)

# Duree : 1788.7s (29.8 min) sur CPU
# C-index train : 0.7756
# C-index test  : 0.7421  (ecart = 0.0001)
\end{codeslide}

\column{0.48\textwidth}
\begin{block}{Hyperparamètres}
\scriptsize
\begin{tabular}{lll}
\toprule
\textbf{Param.} & \textbf{Val.} & \textbf{Rôle} \\
\midrule
\texttt{n\_estimators}     & 50   & nb d'arbres $M$ \\
\texttt{learning\_rate}    & 0,15 & pas $\nu$ \\
\texttt{max\_depth}        & 3    & profondeur \\
\texttt{subsample}         & 0,9  & stochastique \\
\texttt{min\_samples\_leaf}& 10   & régularisation \\
\texttt{loss}              & coxph & log-lik. Cox \\
\bottomrule
\end{tabular}
\end{block}
\vspace{3pt}
\begin{tcolorbox}[mybox,colback=metGreen!10,colframe=metGreen,
  title={C-index de Harrell}]
\scriptsize
\[
\mathrm{C} = \frac{\#\{(i,j)\!:\hat{f}(x_i)>\hat{f}(x_j),T_i<T_j,\delta_i=1\}}
                 {\#\{(i,j)\!:T_i<T_j,\delta_i=1\}}
\]
Train $0{,}776$ / Test $\mathbf{0{,}742}$.
Ecart $= 0{,}0001$ : pas de surapprentissage.
\end{tcolorbox}
\end{columns}
\end{frame}

% ================================================================
% 5. REGLE DE DECISION
% ================================================================
\begin{frame}[fragile,shrink=25]{De $\hat{S}(t\mid x_i)$ à la levée d'alerte}
\vspace{-2pt}
\begin{columns}[T]

\column{0.53\textwidth}
\begin{codeslide}{app/app.py -- calcul scores et horizons}
surv_fns = gbs.predict_survival_function(X)
# Score de confiance : S(30|x_i)
s30 = [float(fn(30)) for fn in surv_fns]

# Horizon T*_i : t tel que S(t) <= S(30)/0.98
t_stars = np.full(len(surv_fns), 30.0)
for i, (fn, s) in enumerate(zip(surv_fns, s30)):
    if s <= 0: continue
    idx = np.searchsorted(-fn.y, -(s/0.98), 'left')
    if idx < len(fn.x):
        t_stars[i] = min(float(fn.x[idx]), 30.0)

# Protocole jury : t_pred = min des horizons confiants
above = df[df['confiance'] > theta]   # theta = 0.30
t_pred = above.groupby('airport_alert_id')['t_hat'].min()
# Gain = sum(t_regle - t_pred) / 3600
\end{codeslide}

\column{0.45\textwidth}
\begin{tcolorbox}[mybox,colback=metBlue!8,colframe=metBlue,
  title={Mécanisme de décision}]
\small
\textbf{1.} $\mathrm{conf}(x_i) = \hat{S}(30 \mid x_i)$\\
\quad proba qu'aucun CG n'arrive dans 30 min.\\[3pt]
\textbf{2.} Si $\mathrm{conf}(x_i) > 0{,}30$ :\\
\quad $\hat{t}_i = \min\{t : \hat{S}(t) \le \hat{S}(30)/0{,}98\}$\\[3pt]
\textbf{3.} $t_{\mathrm{pred}} = \min_i \hat{t}_i$ (conservateur).
\end{tcolorbox}
\vspace{3pt}
\begin{alertblock}{Calibration $\theta$ sur jury}
\footnotesize
Balayage $\theta \in [0;\,0{,}99]$ (100 pas).\\
On maximise le gain sous risque $< 2\,\%$.\\[2pt]
\textbf{$\boldsymbol{\theta = 0{,}30}$ retenu.}\\
Seulement 1\,006/80\,186 strikes confiants.
\end{alertblock}
\end{columns}
\end{frame}

% ================================================================
% 6. RÉSULTATS JURY
% ================================================================
\begin{frame}[shrink=10]{Résultats jury 2023-2025}
\vspace{-4pt}
\begin{columns}[c]
\column{0.25\textwidth}
\begin{tcolorbox}[mybox,colback=metMid!10,colframe=metMid,
  halign=center,top=4pt,bottom=4pt]
{\large\bfseries\color{metMid}1\,352}\\[1pt]
\small alertes évaluées
\end{tcolorbox}
\column{0.25\textwidth}
\begin{tcolorbox}[mybox,colback=metGreen!10,colframe=metGreen,
  halign=center,top=4pt,bottom=4pt]
{\large\bfseries\color{metGreen}103,54 h}\\[1pt]
\small économisées
\end{tcolorbox}
\column{0.25\textwidth}
\begin{tcolorbox}[mybox,colback=metGreen!10,colframe=metGreen,
  halign=center,top=4pt,bottom=4pt]
{\large\bfseries\color{metGreen}0,10\,\%}\\[1pt]
\small risque $<$3 km
\end{tcolorbox}
\column{0.25\textwidth}
\begin{tcolorbox}[mybox,colback=metGreen!10,colframe=metGreen,
  halign=center,top=4pt,bottom=4pt]
{\large\bfseries\color{metGreen}2\,/\,1\,995}\\[1pt]
\small CG manqués
\end{tcolorbox}
\end{columns}

\vspace{3pt}
\begin{columns}[T]
\column{0.52\textwidth}
\begin{block}{Détail par aéroport ($\theta=0{,}30$)}
\scriptsize
\begin{tabular}{lccrr}
\toprule
\textbf{Aéroport} & \textbf{Alertes} & \textbf{Gain (h)} & \textbf{Moy.} & \textbf{Manqués} \\
\midrule
Ajaccio  & 256 & 20,95 & 4,91 min & \textbf{\color{metRed}2} \\
Bastia   & 322 & 25,19 & 4,69 min & 0 \\
Biarritz & 313 & 23,67 & 4,54 min & 0 \\
Nantes   & 122 &  7,89 & 3,88 min & 0 \\
Pise     & 339 & 25,84 & 4,57 min & 0 \\
\bottomrule
\end{tabular}
\end{block}

\column{0.45\textwidth}
\begin{exampleblock}{Conformité sécurité}
\footnotesize
$R = 2/1\,995 = 0{,}10\,\%$ contre limite $2\,\%$ :\\
\textbf{20 fois sous la limite officielle.}
\end{exampleblock}
\vspace{2pt}
\begin{alertblock}{Seule alerte risquée : \texttt{Ajaccio\_835}}
\scriptsize
62 éclairs, 3 dans $<$3 km, 2 manqués après levée.\\
Gain : 62,6 min. Orage avec rebond tardif.\\
Risque local Ajaccio : 1,41\,\%.
\end{alertblock}
\vspace{2pt}
\begin{tcolorbox}[mybox,colback=metBlue!8,colframe=metBlue]
\scriptsize\centering
Modèle actif : \textbf{967/1\,352} alertes (71\,\%).\\
385 alertes restent à la règle 30 min.\\
Gain moyen : \textbf{4,59 min / alerte}.
\end{tcolorbox}
\end{columns}
\end{frame}

% ================================================================
% 7. CONCLUSION
% ================================================================
\begin{frame}{Conclusion}
\vspace{-2pt}
\begin{columns}[T]

\column{0.56\textwidth}
\begin{exampleblock}{Ce que le GBS apporte}
\small
\begin{itemize}\setlength\itemsep{2pt}
  \item[\checkmark] Survie \textbf{individualisée} $\hat{S}(t \mid x_i)$ sans hypothèse paramétrique.
  \item[\checkmark] \textbf{103,54 h économisées} sur jury 2023-2025.
  \item[\checkmark] Risque \textbf{0,10\,\%}, $20\times$ sous les 2\,\%.
  \item[\checkmark] C-index 0,742 -- écart train/test $= 0{,}0001$.
  \item[\checkmark] 1 seule alerte risquée sur 1\,352.
\end{itemize}
\end{exampleblock}
\vspace{3pt}
\begin{alertblock}{Perspectives}
\footnotesize
$\bullet$ Modèle hybride GBS + Weibull AFT (gain + sécurité).\\
$\bullet$ Features v8 : amplitude des éclairs, direction de l'orage.\\
$\bullet$ Déploiement temps réel via l'app Streamlit.
\end{alertblock}

\column{0.42\textwidth}
\begin{tcolorbox}[mybox,colback=metBlue!8,colframe=metBlue,
  title={Pipeline complet}]
\scriptsize
\textbf{1.} \texttt{data\_df\_with\_alert.parquet}\\
\textbf{2.} \texttt{build\_survival()} $\to$ $T_i$ censuré\\
\textbf{3.} \texttt{build\_features\_v7()} $\to$ 13 var.\\
\textbf{4.} \texttt{GradientBoostingSurvivalAnalysis}\\
\quad\texttt{(loss='coxph', n\_estimators=50)}\\
\textbf{5.} \texttt{predict\_survival\_function()}\\
\quad $\hat{S}(t\mid x_i)$ $\to$ $\mathrm{conf} = \hat{S}(30)$\\
\textbf{6.} Sweep $\theta$ $\to$ optimal $= 0{,}30$
\end{tcolorbox}
\vspace{3pt}
\begin{tcolorbox}[mybox,colback=metGreen!12,colframe=metGreen,
  halign=center,top=5pt,bottom=5pt]
\normalsize\bfseries\color{metGreen}
103,54 h économisées\\[2pt]
0,10\,\% de risque\\[1pt]
\small\normalfont\color{black}$20\times$ sous la limite jury
\end{tcolorbox}
\end{columns}
\end{frame}

% ================================================================
% 8. RÉFÉRENCES
% ================================================================
\begin{frame}{Références}
\vspace{-4pt}
\begin{block}{Méthodologie}
\scriptsize
\begin{itemize}\setlength\itemsep{0pt}
  \item Cox, D.R. (1972). \textit{Regression Models and Life-Tables}. J. R. Stat. Soc. B, 34(2).
  \item Breslow, N. (1972). \textit{Contribution to the discussion of Cox (1972)}.
  \item Harrell, F.E. et al. (1982). \textit{Evaluating the yield of medical tests}. JAMA, 247(18).
  \item Friedman, J. (2001). \textit{Greedy Function Approximation: A Gradient Boosting Machine}.
    Annals of Statistics, 29(5), 1189--1232.
  \item Pölsterl, S. (2020). \textit{scikit-survival}. JMLR, 21(212), 1--6.
\end{itemize}
\end{block}
\vspace{2pt}
\begin{columns}[T]
\column{0.50\textwidth}
\begin{tcolorbox}[mybox,colback=metMid!10,colframe=metMid,title={Notebooks}]
\scriptsize
\texttt{modelisation\_v7.ipynb}\\
\texttt{evaluation\_eval\_2023\_v7.ipynb}\\
\texttt{evaluation\_jury\_gbs.ipynb}\\
\texttt{app/app.py} (Streamlit)
\end{tcolorbox}
\column{0.48\textwidth}
\begin{tcolorbox}[mybox,colback=metOrange!10,colframe=metOrange,title={Données}]
\footnotesize
\texttt{data\_df\_with\_alert.parquet}\\
\texttt{segment\_alerts\_all\_airports\_eval.csv}\\
56\,599 éclairs CG, 5 aéroports, 2016-2025.\\
\texttt{models/gbs\_v7\_model.pkl}
\end{tcolorbox}
\end{columns}
\end{frame}

% ── MERCI ───────────────────────────────────────────────────────
{
\setbeamertemplate{footline}{}
\setbeamertemplate{headline}{}
\begin{frame}[plain]
\begin{tikzpicture}[remember picture,overlay]
  \fill[metBlue] (current page.south west) rectangle (current page.north east);
  \fill[metGreen] ([yshift=-1pt]current page.north west)
    rectangle ([yshift=5pt]current page.north east);
  \fill[metMid!50] (current page.south west)
    rectangle ([yshift=2.4cm]current page.south east);
  \node[font=\Huge\bfseries,text=white,anchor=center]
    at ([yshift=0.8cm]current page.center) {Merci};
  \node[font=\large,text=metLight,anchor=center]
    at ([yshift=0.0cm]current page.center) {Questions ?};
  \node[font=\small,text=white!60!metBlue,anchor=center]
    at ([yshift=0.7cm]current page.south)
    {Battle Météorage 2026 · GBS v7 · scikit-survival};
\end{tikzpicture}
\end{frame}
}

\end{document}
